\chapter{构建验证与使用示例}

\section{构建流程}
在安装 TeX Live、XeLaTeX 和 Biber 后，可在项目根目录执行：
\begin{lstlisting}[language=bash,caption={本地构建命令}]
make
make clean
\end{lstlisting}
构建脚本会运行 XeLaTeX、Biber 和必要的重复编译，使目录、引用和参考文献稳定下来。

\section{验证范围}
样例工程覆盖封面、原创性声明、中英文摘要、目录、四个正文级章节、参考文献、附录和致谢。自动检查包括 PDF 是否生成、页数是否非零、关键页眉文字是否出现、中文和英文摘要标题是否存在，以及引用是否能由 Biber 解析。上述检查只能证明样例工程在当前 TeX Live 环境中可构建，不能替代学院对最终论文内容、页数和打印样张的审核。

\section{使用注意事项}
中文摘要应达到学校要求的 400 字以上，英文摘要应达到 350 个实词以上；参考文献应按 GB/T 7714-2015 书写，正文引用顺序与文后编号应保持一致。提交前应在目标打印机或 PDF 阅读器中检查封面、页眉、图表和分页，并根据导师或学院的最新通知更新字段。
